\subsection{Интеграл с переменным верхним (нижним) пределом}
\begin{definition}
	Пусть $\fa\t b\in(a,b)$, $f\in\mc R[a,\t b]$. Тогда $f(\t b)=\int_a^bf(x)dx$ называется \textit{интегралом с переменным верхним пределом}.
\end{definition}
\begin{theorem}\label{9_th1}
	Если $\fa\t b\in(a,b),\ f\in\mc R[a,\t b]$ и $f$ ограничена на $[a,b]$, то $f\in\mc R[a,b]$
\end{theorem}
\begin{proof}
	$f$ - ограничена $\Ra M:=\sup_{[a,b]}|f(x)|,\ (\fa\eps>0)\ \t b:=b-\frac\eps M$\\
	По критерию интегрируемости: $\ex\t P:a=x_0<\dots<x_n=\t b:U(\t P,f)-L(\t P,f)<\frac\eps2$\\
	Введём $P:=\t P\cup\{b\}$.
	\[U(P,f)-L(P,f)=U(\t P,f)-L(\t P,f)+(\sup_{[\t b,b]}f(x)-\inf_{[\t b,b]}f(x))\cdot(b-\t b)<\eps\]
\end{proof}
\begin{note}
	Отметим, что аналогично определяется \textit{интеграл с переменным нижним пределом} и доказывается аналогичная (\ref{9_th1}) для этого интеграла.
\end{note}
\begin{corollary}
	Если $f$ ограничена на $[a,b]$ и имеет лишь конечное число точек разрыва на $[a,b]$, то $f\in\mc R[a,b]$.
\end{corollary}
\begin{proof}
	Возьмём $\{t_i\}_{i=1}^n$, где $t_i$ - точка разрыва $f$.\\
	\[[a,b]=[a,t_1]\cup\left[t_1,\frac{t_1+t_2}2\right]\cup\bigg[\frac{t_1+t_2}2,t_2\bigg]\cup\dots\cup[t_n,b]\]
	Теперь применим ко всем этим подотрезкам $[a,b]$ теорему (\ref{9_th1}).\\
	Применим свойство аддитивности, тогда получим требуемое.
\end{proof}
\begin{note}
	Введём штуку, далее будет применяться:
	\[
	F(x)=
	\begin{cases}
		\int_a^xf(t)dt,\ x>a\\
		0,\ x=a\\
		-\int_b^af(t)dt,\ x<a
	\end{cases}
	\]
\end{note}
\begin{theorem}[Свойства интеграла с переменным верхним пределом]\label{9_th2}
	Если $f\in\mc R[a,b]$, интеграл с переменным верхним пределом $F(x)$ непрерывен на $[a,b]$. Если $f$ непрерывна в $x_0\in[a,b]$, то $F$ дифференцируема в $x_0$, причём $F'(x_0)=f(x_0)$ (если $x_0$ находится на концах отрезка, то односторонне дифференцируема соответственно).
\end{theorem}
\begin{proof}
	$\displaystyle M:=\sup_{[a,b]}|f(x)|,\ \fa x,y:a\le x<y\le b$
	\[|F(y)-F(x)|=\bigg|\int_a^yf(t)dt-\int_a^xf(t)dt\bigg|\le\bigg|\int_x^yf(t)dt\bigg|\le M(y-x)\]
	Получили:
	\[(\fa\eps>0)(\ex\delta=\frac\eps M)(\fa x,y\in[a,b],\ |x-y|<\delta)\ |F(x)-F(y)|\le\eps\]
	Теперь докажем дифференцируемость. Оценим $\displaystyle \bigg|\frac{F(s)-F(x_0)}{s-x_0}-f(x_0)\bigg|$.
	\[\bigg|\frac{F(s)-F(x_0)}{s-x_0}-f(x_0)\bigg|=\bigg|\frac{1}{s-x_0}\int_{x_0}^sf(t)dt-F(x_0)\bigg|=\bigg|\frac{1}{s-x_0}\int_{x_0}^s(f(t)-f(x_0))dt\bigg|\le\eps\]
	$f$ - непрерывна в $x_0\Ra(\fa\eps>0)(\ex\delta>0)(\fa t\in[a,b],\ |t-x_0|\le\delta)\ |f(t)-f(x_0)|<\eps$
\end{proof}
\begin{note}[Лирическое отступление]
	Откуда появилось обозначение интеграла? $\sum$ - обозначение дискретной суммы, тогда как под $\int$ можно понимать суммирование не очень конечных штук, потому он как сумма, только гладкий.
\end{note}
\begin{theorem}[Основная теорема интегрального исчисления]\label{9_th3}
	Если $f\in\mc R[]a,b$ и имеет первообразная $F$ на этом отрезке, то справедлива формула Ньютона - Лейбница:
	\[F(b)-F(a):=\int_a^bf(x)dx=F(x)|_a^b\]
\end{theorem}
\begin{proof}
	$\fa P:a=x_0<\dots<x_n=b$
	\[F(b)-F(a)=(F(b)-x_{n-1})+(F(x_{n-1})-F(x_{n-2}))+\dots+(F(x_1)-F(a))\]
	Применим т. Лагранжа $n$ раз, $t_i\in[x_{i-1},x_i]$:
	\[F(b)-F(a)=F'(t_{n})(b-x_{n-1})+\dots+F'(t_1)(x_1-a)=\sum_{i=1}^nf(t_i)\tr x_i=S(P,f,\{t_i\})\ra\int_a^bf(x)dx\text{, при }\tr(P)\ra0\]
	Значит $\displaystyle F(b)-F(a)=\int_a^bf(x)dx$
\end{proof}
\begin{note}
	Контрпримеры на отсутствие некоторых условий в теореме (\ref{9_th2}) есть, но нынешним инструментарием их не получить.
\end{note}
\begin{note}
	Если $f$ имеет первообразную $F$ на $[a,b]$, то можно определить интеграл по формуле $F(b)-F(a)$.
\end{note}
\begin{corollary}[Формула интегрирования по частям]
	Если $f,f',g,g'\in\mc R[a,b]$, то:
	\[\int_a^bf(x)g'(x)dx=f(b)g(b)-f(a)g(a)-\int_a^bf'(x)g(x)dx\]
\end{corollary}
\begin{theorem}[Формула замены переменной]
	Если $g,g'\in\mc R[\alpha,\beta],\ (\fa r\in[\alpha,\beta])\ m\le g(t)\le M,\ g(\alpha)=a,\ g(\beta)=b$ и $f$ непрерывна на $[m,M]$, то:
	\[\int_a^bf(x)dx=\int_a^bf(g(t))g'(t)dt\]
\end{theorem}
\begin{proof}
	По теореме (\ref{9_th2}): $f$ имеет первообразную $F$.\\
	По теореме (\ref{9_th3}): $\displaystyle\int_a^bf(x)dx=F(b)-F(a)$
	\[F'(x)=f(x),\ (F(g(t)))'=F'(g(t))\cdot g'(t)=f(g(t))\cdot g'(t)\]
	\[\int_\alpha^\beta f(g(t))g'(t)dt=F(g(\beta))-F(g(\alpha))=F(b)-F(a)\]
\end{proof}
\begin{theorem}[Первая теорема о среднем]
	Если $g(x)\ge0,\ \fa x\in[a,b]$, $f,g\in\mc R[a,b]$, $m\le f(x)\le M$, то $\ex\xi\in[m,M]$ такое, что:
	\[\int_a^bf(x)g(x)dx=\xi\int_a^bg(x)dx\]
	Если $f$ непрерывна, то $\displaystyle\ex c\in[a,b]:\int_a^bf(x)g(x)dx=f(c)\int_a^bg(x)dx$
\end{theorem}
\begin{proof}
	$m\le f(x)\le M\Ra m\cdot g(x)\le f(x)\cdot g(x)\le M\cdot g(x)$.
	\[m\int_a^b g(x)dx\le\int_a^b f(x)g(x)dx\le M\int_a^bg(x)dx\]
	\[\xi:=\frac{\int_a^bf(x)g(x)dx}{\int_a^bg(x)dx}\]
	Если $f$ непрерывна, то $f$ принимает все значения на промежутке, значит $\ex c:f(c)=\xi$.\\
	Получили требуемое.
\end{proof}